\chapter{Reproducibilidad}\label{ap:reproducibilidad}

\section{Cómo se repite el trabajo}

El código está en la carpeta \texttt{equipo/benjamin-campos/transfer-learning/} de la rama \texttt{transfer-learning}
del repositorio de NOVA. Estos son los pasos del README. Los tres primeros corren en el PC: preparan los datos y
corren las pruebas. Las redes se entrenan solo en Kaggle, y el rango esperado, que es ridge en CPU, se puede correr en
cualquiera de los dos.

\begin{lstlisting}[basicstyle=\ttfamily\footnotesize]
python -m nova_tl.datos --reposo <carpeta nlsp_tablas> \
       --lifesnaps <daily .csv|.parquet> --salida <carpeta>
python -m pytest tests
kaggle datasets version -p <carpeta con nova_tl/ y los parquet> \
       -m "..." --dir-mode zip
kaggle kernels push -p kaggle --accelerator NvidiaTeslaT4
kaggle kernels output benjacampos/nova-transfer-learning -p <salida>
python -m nova_tl.graficos <salida>/salida   # rehace los gráficos, sin redes
\end{lstlisting}

El entorno del kernel está en la \cref{tab:entorno}. Lo que falta para que una corrida se pueda repetir bit a bit está en
las limitaciones (\cref{sec:limitaciones}): la versión de Chronos no está fija, no se guarda el hash del código y no se
fija el determinismo de CUDA.

\section{Trazabilidad de las cifras}

Las cifras del informe salen de la salida del tercer run del kernel \texttt{benjacampos/nova-transfer-learning}. Para
comprobar que el README describía ese mismo run, se compararon los archivos de la salida del kernel con los que están
en la rama (\cref{tab:traza}): son idénticos. La \cref{tab:hashes} identifica además los datos de entrada y los pesos
publicados en Hugging Face que se usaron para contar los tensores que cambian.

\begin{table}[htbp]
\caption[Archivos del kernel y del repositorio]{Archivos de resultados del tercer run: sha256 de la salida del kernel y
de la copia en la rama (primeros 16 caracteres).}
\label{tab:traza}
\centering\small
\begin{tabular}{lll}\toprule
 Archivo & sha256 (kernel) & sha256 (repositorio) \\
\midrule
\texttt{resultados/metricas.json} & \texttt{ead44910a2e4e2c0} & \texttt{ead44910a2e4e2c0} \\
\texttt{resultados/1\_falsas\_alarmas.png} & \texttt{75e84f37797a317e} & \texttt{75e84f37797a317e} \\
\texttt{resultados/2\_igual\_falsas\_alarmas.png} & \texttt{7c7263a280879b3b} & \texttt{7c7263a280879b3b} \\
\texttt{resultados/3\_curva\_aprendizaje.png} & \texttt{ea81aa0642f760fc} & \texttt{ea81aa0642f760fc} \\
\texttt{resultados/4\_pronostico.png} & \texttt{09e3f31786457416} & \texttt{09e3f31786457416} \\
\bottomrule\end{tabular}

\end{table}

\begin{table}[htbp]
\caption[Datos de entrada y pesos publicados]{Datos de entrada del kernel y pesos de Hugging Face usados en la
auditoría: sha256, primeros 16 caracteres.}
\label{tab:hashes}
\centering\small
\begin{tabular}{ll}\toprule
 Archivo & sha256 (primeros 16) \\
\midrule
\texttt{reposo\_covariables.parquet} & \texttt{736aa6d6c29e225b} \\
\texttt{reposo\_largo.parquet} & \texttt{b20b18b7f6248903} \\
\texttt{tierra\_largo.parquet} & \texttt{db793d9537ba3137} \\
\texttt{nova-tierra} (\texttt{model.safetensors}) & \texttt{3e3ee62a072cdcfa} \\
\texttt{nova-espacio} (\texttt{model.safetensors}) & \texttt{3d1dea0a6435d83c} \\
\bottomrule\end{tabular}

\end{table}

\section{Cómo se construye este informe}

Ninguna cifra de resultado está escrita a mano en los capítulos. Un script lee la salida del kernel y escribe las
macros, tablas y figuras; los capítulos las usan con \verb|\cifra{...}| o las incluyen tal cual. La bibliografía se
armó desde los servicios de metadatos, nunca de memoria. Los pasos son:
\begin{description}[style=nextline, leftmargin=1.2em]
\item[\texttt{resolver\_bib.py}] Busca cada obra en Crossref y DBLP. Acepta un resultado solo si el título
  normalizado coincide (similitud de al menos 0,90), el apellido del primer autor aparece y el año difiere en uno o
  menos. El BibTeX se descarga del propio servicio y solo se le cambia la clave.
\item[\texttt{resolver\_bib2.py}] Segunda pasada, con las mismas reglas, en OpenAlex y arXiv para lo que la primera no
  resolvió. Guarda también los resúmenes de las obras, con los que se contrastó lo que este informe dice de cada una.
\item[\texttt{pulir\_bib.py}] Completa autores y sedes con Semantic Scholar y Crossref, y agrega la palabra clave de
  tema con que se divide la bibliografía.
\item[\texttt{construir\_cifras.py}] Lee la salida del kernel, los datos y los pesos, y escribe
  \texttt{generado/cifras.tex}, las tablas \texttt{tab\_*.tex} y las figuras \texttt{fig\_*.tex}.
\item[\texttt{latexmk}] Compila con Lua\LaTeX{} y biber.
\end{description}

\begin{lstlisting}[basicstyle=\ttfamily\footnotesize]
cd herramientas
G=../generado
python resolver_bib.py $G/referencias_crudas.bib $G/registro_bib.tsv
python resolver_bib2.py $G/referencias_crudas.bib $G/registro_bib2.tsv \
       $G/referencias_resueltas.bib $G/resumenes.json
python pulir_bib.py $G/referencias_resueltas.bib $G/referencias.bib \
       $G/registro_pulir.tsv
python construir_cifras.py .. <carpeta transfer-learning del repositorio>
cd .. && latexmk -lualatex informe
\end{lstlisting}

Tres fuentes no se pudieron resolver con esas reglas y quedaron fuera: el libro \emph{Dataset shift in machine
learning} de Quiñonero-Candela y coautores y el artículo de RevIN de Kim y coautores, que ningún servicio devolvió con
título, autor y año coincidentes, y la página del desafío de Space Apps, que falló la comprobación de su certificado. Los
registros \texttt{registro\_*.tsv} dicen, para cada clave, qué servicio la resolvió y con qué similitud.
